% Body-text fixture for the conference context evaluation.
%
% This is an excerpt from the LaTeX sources of "Attention Is All You Need"
% (Vaswani et al., arXiv:1706.03762), reproduced here ONLY to give the test
% figure a realistic page of running text to sit in. See provenance.md for the
% exact source files, the excerpt boundaries and every edit made.
%
% Nothing on the resulting page is a claim about that paper. The figure this
% file pulls in is a synthetic fixture rendered by this repository from
% tests/data/, and its caption says so.
%
% Venue-neutral on purpose: no \documentclass, no venue macros, no \usepackage.
% The wrappers in tests/conference_context/wrappers/ supply all of that. The one
% thing this file expects from its caller is \input{context-figure}, generated
% per run by the evaluator.

\section{Introduction}

Recurrent neural networks, long short-term memory \citep{hochreiter1997} and gated recurrent \citep{gruEval14} neural networks in particular, have been firmly established as state of the art approaches in sequence modeling and transduction problems such as language modeling and machine translation \citep{sutskever14, bahdanau2014neural, cho2014learning}. Numerous efforts have since continued to push the boundaries of recurrent language models and encoder-decoder architectures \citep{wu2016google,luong2015effective,jozefowicz2016exploring}.

Recurrent models typically factor computation along the symbol positions of the input and output sequences. Aligning the positions to steps in computation time, they generate a sequence of hidden states $h_t$, as a function of the previous hidden state $h_{t-1}$ and the input for position $t$. This inherently sequential nature precludes parallelization within training examples, which becomes critical at longer sequence lengths, as memory constraints limit batching across examples. Recent work has achieved significant improvements in computational efficiency through factorization tricks \citep{Kuchaiev2017Factorization} and conditional computation \citep{shazeer2017outrageously}, while also improving model performance in case of the latter. The fundamental constraint of sequential computation, however, remains.

Attention mechanisms have become an integral part of compelling sequence modeling and transduction models in various tasks, allowing modeling of dependencies without regard to their distance in the input or output sequences \citep{bahdanau2014neural, structuredAttentionNetworks}. In all but a few cases \citep{decomposableAttnModel}, however, such attention mechanisms are used in conjunction with a recurrent network.

In this work we propose the Transformer, a model architecture eschewing recurrence and instead relying entirely on an attention mechanism to draw global dependencies between input and output. The Transformer allows for significantly more parallelization and can reach a new state of the art in translation quality after being trained for as little as twelve hours on eight P100 GPUs.

\section{Background}

The goal of reducing sequential computation also forms the foundation of the Extended Neural GPU \citep{extendedngpu}, ByteNet \citep{NalBytenet2017} and ConvS2S \citep{JonasFaceNet2017}, all of which use convolutional neural networks as basic building block, computing hidden representations in parallel for all input and output positions. In these models, the number of operations required to relate signals from two arbitrary input or output positions grows in the distance between positions, linearly for ConvS2S and logarithmically for ByteNet. This makes it more difficult to learn dependencies between distant positions \citep{hochreiter2001gradient}. In the Transformer this is reduced to a constant number of operations, albeit at the cost of reduced effective resolution due to averaging attention-weighted positions, an effect we counteract with Multi-Head Attention as described in section~\ref{sec:attention}.

% Fixture edit, not part of the source article: start a fresh page here so the
% target figure always lands on an ordinary body page rather than sharing the
% title/abstract page. Everything after this point — margins, running head,
% folio, paragraph density, caption style — is the venue's normal body page.
\clearpage

Self-attention, sometimes called intra-attention is an attention mechanism relating different positions of a single sequence in order to compute a representation of the sequence. Self-attention has been used successfully in a variety of tasks including reading comprehension, abstractive summarization, textual entailment and learning task-independent sentence representations \citep{cheng2016long, decomposableAttnModel, paulus2017deep, lin2017structured}.

End-to-end memory networks are based on a recurrent attention mechanism instead of sequence-aligned recurrence and have been shown to perform well on simple-language question answering and language modeling tasks \citep{sukhbaatar2015}.

% The evaluator writes context-figure.tex into the build directory: the figure
% float carrying the figure under test, its caption, its label, and the TeX-side
% width and page measurements.
\input{context-figure}

To the best of our knowledge, however, the Transformer is the first transduction model relying entirely on self-attention to compute representations of its input and output without using sequence-aligned RNNs or convolution. In the following sections, we will describe the Transformer, motivate self-attention and discuss its advantages over models such as \citep{neural_gpu, NalBytenet2017} and \citep{JonasFaceNet2017}.

\section{Attention}
\label{sec:attention}

An attention function can be described as mapping a query and a set of key-value pairs to an output, where the query, keys, values, and output are all vectors. The output is computed as a weighted sum of the values, where the weight assigned to each value is computed by a compatibility function of the query with the corresponding key.

In practice, we compute the attention function on a set of queries simultaneously, packed together into a matrix $Q$. The keys and values are also packed together into matrices $K$ and $V$. The input consists of queries and keys of dimension $d_k$, and values of dimension $d_v$. We compute the dot products of the query with all keys, divide each by $\sqrt{d_k}$, and apply a softmax function to obtain the weights on the values.
